\section{Reflexion der Methodik: Gütekriterien, Limitationen und Abgrenzungen}
\label{sec:04-04_method-reflection}

Die methodische Qualität der Untersuchung wurde entlang zentraler Gütekriterien qualitativer Forschung reflektiert. Angelehnt an \citeauthor{422:Flick2018} sowie an das Trustworthiness-Konzept nach \citeauthor{423:LincolnGuba1985}werden insbesondere Glaubwürdigkeit (credibility), Übertragbarkeit (transferability), Nachvollziehbarkeit bzw.\ Verlässlichkeit (dependability) und Transparenz bzw.\ Bestätigbarkeit (confirmability) berücksichtigt \cite{422:Flick2018,423:LincolnGuba1985}.

\textbf{Glaubwürdigkeit (Credibility)}: Die Kombination aus \Gls{usability}-Tests, leitfadengestützten Interviews und standardisierten Hintergrunddaten stellt eine methodische Triangulation dar, bei der unterschiedliche Datentypen und Perspektiven auf denselben Gegenstand bezogen werden \cite{422:Flick2018,124:Triangulation2023}. Die standardisierte Durchführung der Test- und Interviewsequenz sowie der einheitliche Leitfaden erhöhen die Vergleichbarkeit der Daten und reduzieren zufällige Einflüsse.

\textbf{Nachvollziehbarkeit und Transparenz (Dependability/Confirmability)}: Alle Schritte der Datenerhebung und -auswertung wurden in Protokollen, Leitfäden und Auswertungsschemata dokumentiert und bilden damit einen nachvollziehbare „audit trail“ im Sinne qualitativer Forschung \cite{422:Flick2018,204:TrustworthinessGuide2020}. Die Vergabe von Interview-IDs (TN-01 bis TN-10) ermöglicht eine klare Zuordnung von Aussagen und Beobachtungen zu einzelnen Fällen, ohne die Anonymität der Teilnehmenden zu gefährden.

\textbf{Übertragbarkeit (Transferability)}: Aufgrund des Einzelfallcharakters und der begrenzten Stichprobe sind die Ergebnisse nicht statistisch generalisierbar. Durch die detaillierte Beschreibung des organisationalen Kontexts, der Nutzergruppen und der Plattformcharakteristika wird jedoch eine „thick description“ bereitgestellt, die es Lesenden erlaubt, eigenständig zu beurteilen, inwieweit die Befunde auf vergleichbare Wiki- oder Intranet-Projekte übertragbar sind \cite{422:Flick2018,123:PracticalGuidance2017}.

\textbf{Limitationen}: Einschränkend wirken die kleine Stichprobengröße sowie der Fokus auf einen frühen Evaluationszeitpunkt im Einführungsprozess. Langfristige Nutzungseffekte, Lerneffekte oder Veränderungen organisationaler Routinen konnten nicht untersucht werden. Zudem basieren die qualitativen Auswertungen überwiegend auf stichpunktartigen Protokollen und nicht auf vollständigen Transkripten, was die Detailtiefe der Analyse begrenzt \cite{421:Mayring2022}. Die quantitative Auswertung bleibt aufgrund der Stichprobengröße auf deskriptive Kennzahlen beschränkt.

\textbf{Abgrenzung}: Gegenstand der Arbeit ist die Wahrnehmung, Nutzbarkeit und Akzeptanz des Wiki-Auftritts sowie deren Zusammenhang mit Informationsarchitektur und \Gls{uxstrategien}. Eine inhaltliche Bewertung fachlicher Dokumente, die Analyse organisationspolitischer Prozesse oder eine systematische Governance-Analyse der gesamten Intranetlandschaft sind explizit nicht Teil der Untersuchung. 

Insgesamt bietet das gewählte Forschungsdesign, trotz der genannten Limitationen, eine fundierte und methodisch begründete Grundlage, um die Forschungsfragen zu beantworten und praxisrelevante Gestaltungsempfehlungen für interne Wissens- und Kollaborationsplattformen abzuleiten \cite{422:Flick2018,423:LincolnGuba1985}.
